\clearpage
\begingroup
\revl
\section{The Klein--Gordon field on Minkowski space}\label{sec:mass}

\noindent\motiv{Chapter~\ref{sec:minkowski} constructs Green data for the wave operator $P$ on
$3+1$ Minkowski space. This chapter adds a mass term. For every real $m^2$ it constructs the
retarded and advanced Green operators of the Klein--Gordon operator $P+m^2$, which gives the free
field with that operator: its CCR net is generating, causal and has the time-slice property for the
inertial slabs, its algebra is primitive given an irreducible representation, its Green operators are
unique, and it is covariant under the orthochronous
Poincar\'e group (Theorem~\ref{thm:kg-field}). At $m^2=0$ the Green data are those of
Chapter~\ref{sec:minkowski} (Corollary~\ref{cor:kg-zero}).}

\medskip\noindent\motiv{The Green operators are the series $G^\pm_{m^2}=\sum_{j\ge0}(-m^2)^j(G^\pm)^{j+1}$ in the
Green operators $G^\pm$ of the wave operator (Definition~\ref{def:kirchhoff}). Formally
$(P+m^2)^{-1}=P^{-1}(1+m^2P^{-1})^{-1}$, and the series is the geometric series of the second factor.
It is a sum of smooth functions $(G^\pm)^{j+1}f$, not of distribution kernels\extraonly{
(Remark~\ref{rem:kg-kernels})}.}

\medskip\noindent\motiv{Already its second term applies $G^+$ to $G^+f$, which is not compactly supported. The
chapter therefore first extends $G^\pm$ to functions supported in the causal future or past of a
compact set and bounds the iterates by $(t-t_0)^{2j}/(2j)!$; these factorials make the series
converge for every real $m^2$ (Proposition~\ref{prop:kg-conv}).\extraonly{ It ends by comparing propagation:
for $m^2=0$ the Green operators propagate only along the light cone, and for every $m^2\neq0$ they
also propagate inside it (Theorem~\ref{thm:kg-iff}).}}

\subsection{The Klein--Gordon operator}

\begin{definition}[Klein--Gordon operator]\label{def:kg}
For a real number $m^2$, the Klein--Gordon operator $P_{m^2}=P+m^2$ acts on $\cF$ and on $\cD$ by
$P_{m^2}h=Ph+m^2h$, with $P$ the wave operator of Definition~\ref{def:mink}. The parameter $m^2$ is
written as a square, as in Remark~\ref{rem:dictionary}, whose operator $P$ is $P_{m^2}$ here, but may
be negative: $m^2>0$ is the free field of mass $m$, and $m^2<0$ is the tachyonic case.
\end{definition}

\begin{lemma}[Properties of the Klein--Gordon operator]\label{lem:kg}
For every $m^2\in\R$ the operator $P_{m^2}$ of Definition~\ref{def:kg} satisfies:
\begin{enumerate}[label=(\alph*),leftmargin=2.5em]
\item $P_0=P$, on $\cF$ and on $\cD$;
\item $\iota P_{m^2}f=P_{m^2}\iota f$ for $f\in\cD$;
\item $\langle P_{m^2}f,h\rangle=\langle f,P_{m^2}h\rangle$ for $f\in\cD$ and $h\in\cF$;
\item $P_{m^2}$ is local, and it is Poincar\'e invariant (Definition~\ref{def:kg-greenops}).
\end{enumerate}
\end{lemma}
\begin{proof}
Each property holds for $P$ (Proposition~\ref{prop:kg-wave}) and for multiplication by the constant
$m^2$, and sums preserve it.
\end{proof}

\subsection{The Green operators beyond test functions}\label{sec:kg-extend}

\noindent\motiv{By Theorem~\ref{thm:kg-greenops}, the axioms, uniqueness and covariance for the field
with $P_{m^2}$ follow once Green operators for $P_{m^2}$ exist; we build them as the series of
Definition~\ref{def:kg-series}, using the extension of $G^\pm$ in Proposition~\ref{prop:kg-ext}.
The value $G^+h(p)$ samples $h$ only on the causal past of $p$. Cutting $h$ off far from that past
turns it into a test function without changing $G^+h$ near $p$, so the results of
Chapter~\ref{sec:minkowski} for test functions carry over.}

\begin{definition}[Extended Green operators; functions supported in a causal cone]\label{def:kg-supported}
A function $h:M\to\R$ is \emph{supported in the future of a compact set} if $\supp h\subseteq J^+(K)$
for some compact $K\subseteq M$, and \emph{supported in the past of a compact set} if
$\supp h\subseteq J^-(K)$ for some compact $K$. $\cF^+$ and $\cF^-$ are the smooth such functions;
test functions lie in both.

The \emph{extended Green operators} are the Green operators $G^\pm$ of
Definition~\ref{def:kirchhoff}, extended to all functions: for $h:M\to\R$ and $p\in M$,
\revm{$G^\pm h(p)=\int_Mh(p-w)\,\mu^\pm(dw)$,} \artefact{taken as $0$ when it does not converge
absolutely}\revm{. For continuous $h$ this is
\[
G^\pm h(p)=\frac1{4\pi}\int_{\R^3}h(p^0\mp|y|,\vec p-y)\,\frac{dy}{|y|}
\]
by the proof of Proposition~\ref{prop:kirchhoff}(a).} On test functions they are the operators of Definition~\ref{def:kirchhoff}, by
Proposition~\ref{prop:kirchhoff}(a). Their iterates are $(G^\pm)^0h=h$ and
$(G^\pm)^{j+1}h=G^\pm\bigl((G^\pm)^jh\bigr)$.
\end{definition}

\begin{lemma}[Causality of $G^\pm$]\label{lem:kg-causal}
Let $h,h':M\to\R$ be continuous and $p\in M$. If $h=h'$ on $J^-(\{p\})$, then the extended Green operator
$G^+$ of Definition~\ref{def:kg-supported} satisfies $G^+h(p)=G^+h'(p)$. If $h=h'$ on $J^+(\{p\})$,
then $G^-h(p)=G^-h'(p)$.
\end{lemma}
\begin{proof}
The sample point $q=(p^0-|y|,\vec p-y)$ of $G^+h(p)$ has $p-q=(|y|,y)$ future causal, so $q\le p$,
and the two integrands agree for every $y$. The sample points of $G^-h(p)$ are
$(p^0+|y|,\vec p-y)\ge p$.
\end{proof}

\begin{lemma}[\revm{Localization}]\label{lem:kg-cutoff}
Let $M$, $\le$, $J^\pm$ and $\cD$ be as in Definition~\ref{def:mink}, and $\operatorname{dist}$ the
distance of the norm $\max(|q^0|,|\vec q|)$ on $M$.
\begin{enumerate}[label=(\alph*),leftmargin=2.5em]
\item For compact $K,K'\subseteq M$, the set $J^+(K)\cap J^-(K')$ is compact.
\item For compact $A\subseteq M$ there is $\psi\in\cD$ with $\psi=1$ on
$\{q:\operatorname{dist}(q,A)<1\}$.
\item Let $h,\psi:M\to\R$ and $K,U\subseteq M$ with $\supp h\subseteq J^+(K)$ and $\psi=1$ on
$\{q:\operatorname{dist}(q,A)<1\}$, where $A=J^+(K)\cap J^-(\overline U)$. Then for every $p\in U$, $\psi h=h$ on a
neighbourhood of every point of $J^-(\{p\})$.
\end{enumerate}
\end{lemma}
\begin{proof}
\begin{enumerate}[leftmargin=2.5em]
\item[(a)] The set is closed (Lemma~\ref{lem:mink-compact}(b)). It lies in
$J^+(K)\cap\{q^0\le t_2\}$, with $t_2$ \revm{an upper bound of the times in} $K'$, which is compact by
Lemma~\ref{lem:mink-compact}(b) for the future unit vector $e_0$.
\item[(b)] Take a smooth bump function equal to $1$ on the closed $1$-neighbourhood of $A$ and supported in
its $2$-neighbourhood, which is compact.
\item[(c)] Let $q\le p$. If $q\in\supp h$, then $q\in A$ and $\psi=1$ near $q$. Otherwise $h=0$ near
$q$, and so is $\psi h$.
\end{enumerate}
\end{proof}

\begin{proposition}[$G^\pm$ on $\cF^\pm$]\label{prop:kg-ext}
For the extended Green operators $G^\pm$ and the spaces $\cF^\pm$ of Definition~\ref{def:kg-supported}:
\begin{enumerate}[label=(\alph*),leftmargin=2.5em]
\item \revm{if $h$ is continuous with $\supp h\subseteq J^+(K)$, $K$ compact, then
$\supp G^+h\subseteq J^+(K)$;}
\item \revm{for $h\in\cF^+$,} $G^+h$ is smooth, and $\partial_vG^+h=G^+\partial_vh$ for every direction $v$;
\item \revm{for $h\in\cF^+$,} $G^+h\in\cF^+$;
\item \revm{for $h\in\cF^+$,} $PG^+h=h$ and $G^+Ph=h$;
\item \revm{with $r(t,\vec x)=(-t,\vec x)$, $G^-h(p)=G^+(h\circ r)(rp)$ for continuous $h$, and
(a)--(d) hold for $G^-$ on $\cF^-$ with $J^-$ in place of $J^+$.}
\end{enumerate}
\end{proposition}
\begin{proof}
\emph{Supports.} If $G^+h(p)\ne0$, some sample point $q\le p$ has $h(q)\ne0$, so $p\in J^+(K)$ by
transitivity; $J^+(K)$ is closed. This gives \revm{(a)}\revo{, and with (b) below, (c)}.

\emph{Localizing.} Fix $p'$, let $U$ be the open unit ball about it in the norm of
Lemma~\ref{lem:kg-cutoff}, and take $\psi$ from Lemma~\ref{lem:kg-cutoff}(b), (a) with
$A=J^+(K)\cap J^-(\overline U)$. For $p\in U$, $h=\psi h$ on $J^-(\{p\})$ by
Lemma~\ref{lem:kg-cutoff}(c), so $G^+h=G^+(\psi h)$ on $U$ by Lemma~\ref{lem:kg-causal}. Since $\psi h\in\cD$,
Proposition~\ref{prop:kirchhoff}(b) gives smoothness on $U$. For the derivative,
$\partial_vh=\partial_v(\psi h)$ near every point of $J^-(\{p\})$ by Lemma~\ref{lem:kg-cutoff}(c),
so $\partial_vG^+h=G^+\partial_v(\psi h)=G^+\partial_vh$ on $U$. This is \revm{(b)}.

\emph{Green identities.} On $U$, $PG^+h=PG^+(\psi h)=\psi h$ by locality of $P$ and
Corollary~\ref{cor:mink-green}, and $\psi h(p)=h(p)$. Since $Ph=P(\psi h)$ near every point of
$J^-(\{p\})$, also $G^+Ph(p)=G^+P(\psi h)(p)=h(p)$. This is \revm{(d)}.

\emph{The advanced case.} $\mu^-$ is the image of $\mu^+$ under $r$, since $dy/|y|$ is invariant
under $y\mapsto-y$; this gives the formula for $G^-$. The involution $r$ reverses the order, so it
exchanges $\cF^+$ and $\cF^-$; it commutes with $P$, and
$\partial_v(h'\circ r)=(\partial_{rv}h')\circ r$ for smooth $h'$. Applying the retarded statements to
$h\circ r$ gives the advanced ones\revm{, (e)}.
\end{proof}

\begin{lemma}[Power bound]\label{lem:kg-power}
For every $s\ge0$ and every integer $n\ge0$,
\[
\int_0^s\rho(s-\rho)^n\,d\rho=\frac{s^{n+2}}{(n+1)(n+2)}.
\]
Let $h:M\to\R$ be continuous, $n\ge0$ an integer, and $t_0,T,C\in\R$, with $h(q)=0$ for $q^0<t_0$
and $|h(q)|\le C(q^0-t_0)^n/n!$ for $t_0\le q^0\le T$. Then the extended Green operator $G^+$ of
Definition~\ref{def:kg-supported} satisfies
\[
|G^+h(p)|\le C\,\frac{(p^0-t_0)^{n+2}}{(n+2)!}\qquad\text{for }t_0\le p^0\le T.
\]
\end{lemma}
\begin{proof}
\emph{The integral.} Substitute $\rho\mapsto s-\rho$ and integrate $(s-\rho)\rho^n$ term by term.

\emph{The bound.} Put $s=p^0-t_0$. The sample point $(p^0-|y|,\vec p-y)$ has time at most $T$, and
$h$ vanishes there unless $|y|\le s$, when its time exceeds $t_0$ by $s-|y|$. In polar coordinates,
\[
|G^+h(p)|\le\frac C{4\pi}\int_{|y|\le s}\frac{(s-|y|)^n}{n!\,|y|}\,dy
=\frac C{n!}\int_0^s\rho(s-\rho)^n\,d\rho=C\,\frac{s^{n+2}}{(n+2)!}.
\]
\end{proof}

\begin{corollary}[Iterates]\label{cor:kg-iterate}
Let $h$ be smooth with $\supp h\subseteq J^+(K)$, $K$ compact, let $j\ge0$, and let $(G^+)^j$ be the
iterate of Definition~\ref{def:kg-supported}. Then:
\begin{enumerate}[label=(\alph*),leftmargin=2.5em]
\item $(G^+)^jh$ is smooth with support in $J^+(K)$, and for every $t_0\in\R$ it vanishes on
$\{q^0<t_0\}$ when $h$ does;
\item if $t_0,T,C\in\R$, $h$ vanishes on $\{q^0<t_0\}$ and $|h|\le C$ on $\{q^0\le T\}$, then
$|(G^+)^jh(p)|\le C(p^0-t_0)^{2j}/(2j)!$ for $t_0\le p^0\le T$;
\item $\partial_v(G^+)^jh=(G^+)^j\partial_vh$ for every direction $v$; so if $t_0,T,C'\in\R$, $h$
vanishes on $\{q^0<t_0\}$ and $|\partial_vh|\le C'$ on $\{q^0\le T\}$, then
$|\partial_v(G^+)^jh(p)|\le C'(p^0-t_0)^{2j}/(2j)!$ for $t_0\le p^0\le T$.
\end{enumerate}
\end{corollary}
\begin{proof}
Induction on $j$.
\begin{enumerate}[leftmargin=2.5em]
\item[(a)] Proposition~\ref{prop:kg-ext}\revm{(b), (c)}; the vanishing before $t_0$ follows from
Lemma~\ref{lem:kg-causal}, since all sample points of $G^+h(p)$ lie at times at most $p^0$.
\item[(b)] Lemma~\ref{lem:kg-power} with $n=2j$, since $2j+2=2(j+1)$.
\item[(c)] Proposition~\ref{prop:kg-ext}\revm{(b)}, applied $j$ times, and (b) for $\partial_vh$, which
vanishes on the open set $\{q^0<t_0\}$.
\end{enumerate}
\end{proof}

\begin{proposition}[Formal adjointness of $G^\pm$ on $\cF^\pm$]\label{prop:kg-adj}
For $h_1\in\cF^-$ and $h_2\in\cF^+$ (Definition~\ref{def:kg-supported}), the functions $h_1\,G^+h_2$
and $(G^-h_1)\,h_2$ are integrable, and $\int_Mh_1\,G^+h_2\,dp=\int_M(G^-h_1)\,h_2\,dp$.
\end{proposition}
\begin{proof}
\emph{Integrability.} Let $\supp h_1\subseteq J^-(K')$ and $\supp h_2\subseteq J^+(K)$. By
Proposition~\ref{prop:kg-ext} both products vanish outside $A=J^+(K)\cap J^-(K')$, which is compact
(Lemma~\ref{lem:kg-cutoff}(a)); they are continuous, hence integrable.

\emph{Cutting off.} Take $\psi\in\cD$ equal to $1$ near $A$. For $p\in A$ and $q\le p$, either
$h_2(q)=0$ or $q\in J^+(K)\cap J^-(K')=A$; so $\psi h_2=h_2$ on $J^-(\{p\})$, and likewise $\psi
h_1=h_1$ on $J^+(\{p\})$. Hence $G^+h_2(p)=G^+(\psi h_2)(p)$ and $G^-h_1(p)=G^-(\psi h_1)(p)$ by
Lemma~\ref{lem:kg-causal}, and the \revm{localized} products also vanish outside $A$. So
Proposition~\ref{prop:kirchhoff}(d) for the test functions $\psi h_1$, $\psi h_2$ gives
$\int h_1\,G^+h_2=\int\psi h_1\,G^+(\psi h_2)=\int G^-(\psi h_1)\,\psi h_2=\int(G^-h_1)\,h_2$.
\end{proof}

\subsection{The Green operators as a series}

\begin{definition}[The series]\label{def:kg-series}
For $m^2\in\R$ and $f:M\to\R$, with the iterates $(G^\pm)^{j+1}$ of the extended Green operators
(Definition~\ref{def:kg-supported}),
\[
G^\pm_{m^2}f(p)=\sum_{j\ge0}(-m^2)^j(G^\pm)^{j+1}f(p),
\]
with $(-m^2)^0=1$, the sum taken pointwise (\artefact{and as $0$ where the series does not converge
absolutely}). For test functions the series converges absolutely (Proposition~\ref{prop:kg-conv}), so
the convention is not used for them.
\end{definition}

\begin{lemma}[Linearity of the iterates]\label{lem:kg-linear}
For the extended Green operators $G^\pm$ and their iterates (Definition~\ref{def:kg-supported}): for
continuous $h$ supported in the future of a compact set, the integrand $y\mapsto
h(p^0-|y|,\vec p-y)/|y|$ of $G^+h(p)$ is integrable for every $p$, and $G^+$ is additive on such
functions. For smooth $h_1,h_2$ supported in the future of compact sets and every $j$,
$(G^+)^j(h_1+h_2)=(G^+)^jh_1+(G^+)^jh_2$. For every $h:M\to\R$ and $\lambda\in\R$,
$(G^\pm)^j(\lambda h)=\lambda\,(G^\pm)^jh$ and $(G^\pm)^{j+1}h=(G^\pm)^j(G^\pm h)$.
\end{lemma}
\begin{proof}
\emph{Integrability and additivity.} Let $\supp h\subseteq J^+(K)$; if $K$ is empty, $h=0$.
Otherwise let $t_K=\min_Kq^0$. A sample
point in $J^+(K)$ has time at least $t_K$, so the integrand vanishes unless $|y|\le p^0-t_K$. There it
is bounded by $\max_B|h|/|y|$, with $B$ the compact set of the sample points $(p^0-|y|,\vec p-y)$,
$|y|\le p^0-t_K$, and $1/|y|$ is integrable on balls of $\R^3$. Additivity of $G^+$ follows.

\emph{Iterates.} Additivity of $(G^+)^j$ follows by induction with Proposition~\ref{prop:kg-ext}\revm{(b),
(c)}. Homogeneity holds for every integral, and the last identity is associativity of iteration.
\end{proof}

\begin{proposition}[Convergence]\label{prop:kg-conv}
Let $m^2,t_0\in\R$, $C\ge0$, and let $f\in\cD$ vanish on $\{q^0<t_0\}$ with $|f|\le C$. \revm{For the
iterates of Definition~\ref{def:kg-supported} and the series of Definition~\ref{def:kg-series}:}
\begin{enumerate}[label=(\alph*),leftmargin=2.5em]
\item \revm{for} every $j\ge0$, $p\in M$ and $T\ge t_0$ with $p^0\le T$,
\[
\bigl|(-m^2)^j(G^+)^{j+1}f(p)\bigr|\le C\,(T-t_0)^2\,\frac{\bigl(\sqrt{|m^2|}\,(T-t_0)\bigr)^{2j}}{(2j)!}.
\]
\item \revm{the series $G^+_{m^2}f(p)$ converges absolutely at every $p$, and on each set
$\{q^0\le T\}$ its terms are bounded by a summable sequence that does not depend on $p$; since
every test function vanishes before some time and is bounded, this holds for every test function;}
\item \revm{$G^+_{m^2}(f+g)=G^+_{m^2}f+G^+_{m^2}g$ for $f,g\in\cD$, and $G^+_{m^2}(\lambda f)=\lambda\,G^+_{m^2}f$
for every $f:M\to\R$ and $\lambda\in\R$.}
\end{enumerate}
\end{proposition}
\begin{proof}
\emph{The bound.} Corollary~\ref{cor:kg-iterate}(b) with $j+1$ in place of $j$ gives
$C(p^0-t_0)^{2j+2}/(2j+2)!$, which is at most the right side, since $(2j+2)!\ge(2j)!$ and
$p^0-t_0\le T-t_0$; for $p^0<t_0$ the term vanishes (Corollary~\ref{cor:kg-iterate}(a)).

\emph{Convergence.} On $\{q^0\le T\}$ the bound is the $j$-th term of $C(T-t_0)^2\cosh\xi$ with
$\xi=\sqrt{|m^2|}(T-t_0)$\revm{, a convergent series, since $\sum_j\xi^{2j}/(2j)!=\cosh\xi$ for every
real $\xi$}.

\emph{Linearity.} Each term is additive on $\cD$ (Lemma~\ref{lem:kg-linear}), and absolutely
convergent series add termwise. Each term is homogeneous.
\end{proof}

\begin{proposition}[Derivatives and smoothness]\label{prop:kg-smooth}
Let $e_0,\dots,e_3$ be the standard basis of $\R^4$ and $e^0,\dots,e^3$ its dual basis. For every
$f\in\cD$ and $m^2\in\R$, the series $G^+_{m^2}f$ of Definition~\ref{def:kg-series} satisfies:
\begin{enumerate}[label=(\alph*),leftmargin=2.5em]
\item $G^+_{m^2}f$ is differentiable at every $p\in M$, with derivative
$\sum_{i=0}^3G^+_{m^2}(\partial_{e_i}f)(p)\,e^i$;
\item $\partial_vG^+_{m^2}f=G^+_{m^2}\partial_vf$ for every $v\in\R^4$;
\item $G^+_{m^2}f$ is smooth.
\end{enumerate}
\end{proposition}
\begin{proof}
\begin{enumerate}[leftmargin=2.5em]
\item[(a)] Fix $p$ and the open unit ball $U$ about it. Each term $H_j=(-m^2)^j(G^+)^{j+1}f$ is
smooth with $\partial_{e_i}H_j=(-m^2)^j(G^+)^{j+1}\partial_{e_i}f$ (Corollary~\ref{cor:kg-iterate}). By
Proposition~\ref{prop:kg-conv} for the test functions $\partial_{e_i}f$, with $t_0$ below $\supp f$
and $T=\max(t_0,p^0+1)$, the norms of the derivatives of $H_j$ are bounded on $U$ by a summable
sequence in $j$.

A series of differentiable functions that converges at one point of a convex open set, with
derivatives bounded there by a summable sequence, is differentiable on the set with the summed
derivative. Summing gives the formula.
\item[(b)] Write $v=\sum_iv^ie_i$ and use (a), additivity and homogeneity
(Proposition~\ref{prop:kg-conv}).
\item[(c)] By induction on $n$, $G^+_{m^2}f$ is $n$ times continuously differentiable for every test
function $f$: by (a) its derivative is $p\mapsto\sum_iG^+_{m^2}(\partial_{e_i}f)(p)\,e^i$, and each
$\partial_{e_i}f$ is a test function.
\end{enumerate}
\end{proof}

\begin{theorem}[Green identities]\label{thm:kg-green-id}
For every $f\in\cD$ and $m^2\in\R$, the series $G^+_{m^2}$ of Definition~\ref{def:kg-series} and the
operator $P_{m^2}$ of Definition~\ref{def:kg} satisfy $P_{m^2}G^+_{m^2}f=f$ and $G^+_{m^2}P_{m^2}f=f$.
\end{theorem}
\begin{proof}
\emph{Commuting with $P$.} $P$ is a combination of second directional derivatives, so
Proposition~\ref{prop:kg-smooth}(b) twice, additivity and homogeneity give $PG^+_{m^2}f=G^+_{m^2}Pf$.

\emph{Telescoping.} By Lemma~\ref{lem:kg-linear} and Corollary~\ref{cor:mink-green},
$(G^+)^{j+1}Pf=(G^+)^j(G^+Pf)=(G^+)^jf$. Hence, shifting the index of an absolutely convergent
series,
\[
G^+_{m^2}Pf=\sum_{j\ge0}(-m^2)^j(G^+)^jf=f+\sum_{j\ge0}(-m^2)^{j+1}(G^+)^{j+1}f=f-m^2G^+_{m^2}f.
\]
So $P_{m^2}G^+_{m^2}f=G^+_{m^2}Pf+m^2G^+_{m^2}f=f$, and $G^+_{m^2}P_{m^2}f=f$ by additivity and
homogeneity.
\end{proof}

\begin{proposition}[Supports]\label{prop:kg-supp}
For every $f\in\cD$ and $m^2\in\R$, the series of Definition~\ref{def:kg-series} satisfies
$\supp G^+_{m^2}f\subseteq J^+(\supp f)$.
\end{proposition}
\begin{proof}
Every term vanishes outside $J^+(\supp f)$ (Corollary~\ref{cor:kg-iterate}(a)), so the sum does, and
$J^+(\supp f)$ is closed.
\end{proof}

\begin{proposition}[The advanced series]\label{prop:kg-adv}
For every $f\in\cD$, $m^2\in\R$ and $j\ge0$, with $r(t,\vec x)=(-t,\vec x)$ and the iterates and series
of Definitions~\ref{def:kg-supported} and~\ref{def:kg-series},
$(G^-)^jf=\bigl((G^+)^j(f\circ r)\bigr)\circ r$ and $G^-_{m^2}f=\bigl(G^+_{m^2}(f\circ r)\bigr)\circ r$.
\revm{Consequently:}
\begin{enumerate}[label=(\alph*),leftmargin=2.5em]
\item the series $G^-_{m^2}f(p)$ converges absolutely at every $p$;
\item $G^-_{m^2}$ is additive on $\cD$ and homogeneous;
\item $\partial_vG^-_{m^2}f=G^-_{m^2}\partial_vf$;
\item $G^-_{m^2}f$ is smooth;
\item $P_{m^2}G^-_{m^2}f=f$ and $G^-_{m^2}P_{m^2}f=f$;
\item $\supp G^-_{m^2}f\subseteq J^-(\supp f)$.
\end{enumerate}
\end{proposition}
\begin{proof}
\emph{The identities.} The one-step identity is Proposition~\ref{prop:kg-ext}. By induction, applying
it to $h=((G^+)^j(f\circ r))\circ r$ gives $(G^-)^{j+1}f=((G^+)^{j+1}(f\circ r))\circ r$, since
$r\circ r$ is the identity. Summing gives the series identity.

\emph{Consequences.} The reflection commutes with $P$ and exchanges $J^+$ and $J^-$, so the
consequences follow from Propositions~\ref{prop:kg-conv}, \ref{prop:kg-smooth} and~\ref{prop:kg-supp}
and Theorem~\ref{thm:kg-green-id} for $f\circ r$.
\end{proof}

\begin{proposition}[Formal adjointness of the series]\label{prop:kg-series-adj}
For $f,g\in\cD$ and $m^2\in\R$, the series of Definition~\ref{def:kg-series} satisfy
$\int_Mf\,G^+_{m^2}g\,dp=\int_M(G^-_{m^2}f)\,g\,dp$.
\end{proposition}
\begin{proof}
\emph{Exchanging sum and integral.} On the compact set $\supp f$, the terms of $G^+_{m^2}g$ are
bounded by a summable sequence: Proposition~\ref{prop:kg-conv} for $g$, with $t_0$ below $\supp g$ and
\revm{$T=\max(t_0,T')$ with $T'$ an upper bound of the times in $\supp f$}. So $\int f\,G^+_{m^2}g=\sum_j(-m^2)^j\int f\,(G^+)^{j+1}g$ by dominated
convergence. Likewise on the other side, where the terms of $G^-_{m^2}f$ are bounded on $\supp g$
through $(G^-)^{j+1}f=((G^+)^{j+1}(f\circ r))\circ r$ (Proposition~\ref{prop:kg-adv}) and
Proposition~\ref{prop:kg-conv} for $f\circ r$.

\emph{Moving one $G$ at a time.} For $0\le i\le j$, Proposition~\ref{prop:kg-adj} with
$h_1=(G^-)^if\in\cF^-$ and $h_2=(G^+)^{j-i}g\in\cF^+$ gives
$\int(G^-)^if\,(G^+)^{j+1-i}g=\int(G^-)^{i+1}f\,(G^+)^{j-i}g$. Here $h_1$ is supported in the past of
$\supp f$ by the advanced case of Proposition~\ref{prop:kg-ext}, applied $i$ times, and $h_2$ in the
future of $\supp g$ by Corollary~\ref{cor:kg-iterate}(a). The $j+1$ steps give
$\int f\,(G^+)^{j+1}g=\int(G^-)^{j+1}f\,g$; summing gives the claim.
\end{proof}

\subsection{\revm{The Klein--Gordon field}}

\begin{theorem}[\revm{The Klein--Gordon field}]\label{thm:kg-field}
For every $m^2\in\R$, the series $G^\pm_{m^2}$ of Definition~\ref{def:kg-series}, restricted to test
functions, are Green operators (Definition~\ref{def:kg-greenops}) for the Klein--Gordon operator
$P_{m^2}$ of Definition~\ref{def:kg}. Consequently:
\begin{enumerate}[label=(\alph*),leftmargin=2.5em]
\item they give Green data with the inertial slabs, whose CCR net is generating, causal and has the
time-slice property over the derived causal index; if the CCR algebra has an irreducible
representation, its algebra is moreover primitive;
\item every map $G':\cD\to\cF$ with $P_{m^2}G'f=f$ and $\supp G'f\subseteq J^+(\supp f)$ for all
$f\in\cD$ equals $G^+_{m^2}$, and likewise $G^-_{m^2}$ with $J^-$;
\item the CCR net is covariant under the orthochronous Poincar\'e group $\mathcal P^\uparrow$, with
$\kappa_\gamma(O)=\gamma(O)$ and $\alpha_\gamma$ induced by the symmetry of
Theorem~\ref{thm:kg-greenops}(c).
\end{enumerate}
\end{theorem}
\begin{proof}
$G^\pm_{m^2}$ map $\cD$ linearly into $\cF$ (Propositions~\ref{prop:kg-conv}, \ref{prop:kg-smooth}
and~\ref{prop:kg-adv}). Condition (a) of Definition~\ref{def:kg-greenops} is
Theorem~\ref{thm:kg-green-id} and Proposition~\ref{prop:kg-adv}, (b) is
Propositions~\ref{prop:kg-supp} and~\ref{prop:kg-adv}, and (c) is
Proposition~\ref{prop:kg-series-adj}; $P_{m^2}$ has the required properties by Lemma~\ref{lem:kg}.
Then Theorem~\ref{thm:kg-greenops}; generation, causality and time slice hold without the hypothesis,
as after Theorem~\ref{thm:green-axioms}.
\end{proof}
\noindent\motiv{The cutoffs, and with them the derived causal index, do not depend on $m^2$. So
Remark~\ref{rem:mink-slabs} applies unchanged.}

\begin{corollary}[The massless field is the case $m^2=0$]\label{cor:kg-zero}
The series of Definition~\ref{def:kg-series} satisfy $G^\pm_0f=G^\pm f$ for every $f:M\to\R$, and the
Green data of Theorem~\ref{thm:kg-field} with $m^2=0$ are equal to the Green data of the free massless
field (Theorem~\ref{thm:mink-green}). So every statement of Chapter~\ref{sec:minkowski} about the
massless field is a statement about $P_0$.
\end{corollary}
\begin{proof}
All terms with $j\ge1$ carry the factor $0^j=0$. By Lemma~\ref{lem:kg}(a), the Green data of
Theorem~\ref{thm:kg-field} at $m^2=0$ are those of Theorem~\ref{thm:kg-greenops} for $P$ and the
extended Green operators, which are the Green data of Theorem~\ref{thm:mink-green} by
Proposition~\ref{prop:kg-wave}. Green data are determined by their components, since (G1)--(G7) are properties, so equal
components give equal Green data.
\end{proof}

\begin{extra}
\subsection{Huygens' principle}

\noindent\motiv{For $m^2=0$ the operator $G^+$ samples $f$ only on the past light-cone surface of
the point, so $G^+f$ vanishes at every point strictly inside the future cone of all of $\supp f$:
signals travel exactly at the speed of light (Proposition~\ref{prop:kg-huygens}). For $m^2\neq0$ the
second term $-m^2(G^+)^2f$ already reaches inside the cone, and Theorem~\ref{thm:kg-fail} shows that,
for suitable $f$ and $p$, the rest of the series does not cancel it.}

\begin{definition}[Light-cone surface, chronological order\revm{, Huygens' principle}]\label{def:kg-cone}
$C^+=\{v\in M:v^0\ge|\vec v|\}$ is the closed future cone, and $\partial C^+=\{v:v^0=|\vec v|\}$ is
its boundary, the future light-cone surface. For $p,q\in M$, $p\ll q$ means $q-p\in I^+$, that is
$|\vec q-\vec p|<q^0-p^0$ ($I^+$ as in Definition~\ref{def:mink}). For $A,B\subseteq M$,
$A+B=\{a+b:a\in A,\ b\in B\}$ and $-B=\{-b:b\in B\}$; so $\supp f+\partial C^+$ is the union of the
light-cone surfaces $q+\partial C^+$, $q\in\supp f$. \revm{A \revo{linear} map $G:\cD\to\cF$ with the supports of
$G^+$ in Definition~\ref{def:kg-greenops}(b) satisfies \emph{Huygens' principle} if $Gf(p)=0$ for every
$f\in\cD$ and every $p\in M$ with $q\ll p$ for all $q\in\supp f$; one with the supports of $G^-$
satisfies it if $Gf(p)=0$ whenever $p\ll q$ for all $q\in\supp f$.}
\end{definition}

\begin{proposition}[Huygens' principle for the wave operator]\label{prop:kg-huygens}
For every $f\in\cD$, the extended Green operators of Definition~\ref{def:kg-supported} and the sets of
Definition~\ref{def:kg-cone} satisfy $\supp G^+f\subseteq\supp f+\partial C^+$ and
$\supp G^-f\subseteq\supp f+(-\partial C^+)$. In particular $G^+f(p)=0$ whenever $q\ll p$ for all
$q\in\supp f$, and $G^-f(p)=0$ whenever $p\ll q$ for all $q\in\supp f$.
\end{proposition}
\begin{proof}
\emph{Supports.} Every sample point of $G^+f(p)$ is $p-(|y|,y)$ with $(|y|,y)\in\partial C^+$. If
$p\notin\supp f+\partial C^+$, no sample point lies in $\supp f$, and $G^+f(p)=0$. The set
$\supp f+\partial C^+$ is closed, as the sum of a compact and a closed set, so it contains
$\supp G^+f$. For $G^-$ the sample points are $p-(-|y|,y)$ with $(-|y|,y)\in-\partial C^+$.

\emph{Points inside the cone.} If $q\ll p$ for all $q\in\supp f$, then $p-q\in I^+$, which misses
$\partial C^+$, so $p\notin\supp f+\partial C^+$. The advanced case is the same with $-I^+$.
\end{proof}

\begin{lemma}[Positivity of $G^+$]\label{lem:kg-pos}
Let $h:M\to\R$ be continuous and nonnegative, $p\in M$, and let $y\mapsto h(p^0-|y|,\vec p-y)/|y|$ be
integrable on $\R^3$. If $h(p^0-|y_1|,\vec p-y_1)>0$ for some $y_1\ne0$, then the extended Green operator
of Definition~\ref{def:kg-supported} satisfies $G^+h(p)>0$.
\end{lemma}
\begin{proof}
The integrand is nonnegative, and continuous and positive at $y_1\neq0$, hence positive on a ball
about $y_1$, which has positive measure.
\end{proof}

\begin{proposition}[Positivity of $(G^+)^2$]\label{prop:kg-pos2}
Let $f\in\cD$ be nonnegative with $f(q_0)>0$, and let $q_0\ll p$ (Definition~\ref{def:kg-cone}). Then
the iterate of Definition~\ref{def:kg-supported} satisfies $(G^+)^2f(p)>0$.
\end{proposition}
\begin{proof}
\emph{A point on both cones.} Put $s=p^0-q_0^0$ and $d=\vec p-\vec q_0$, so $|d|<s$. Choose a
unit vector $\hat d$ with $d=|d|\hat d$ (any unit vector if $d=0$) and let
$y_1=\tfrac12(s+|d|)\hat d\ne0$. The sample point $q=(p^0-|y_1|,\vec p-y_1)$ has
$q^0-q_0^0=\tfrac12(s-|d|)>0$ and $\vec q-\vec q_0=\tfrac12(|d|-s)\hat d$. So $z_1=\vec q-\vec q_0$
satisfies $|z_1|=q^0-q_0^0>0$, and $(q^0-|z_1|,\vec q-z_1)=q_0$.

\emph{Two applications of Lemma~\ref{lem:kg-pos}.} For $h=f$ at $q$ with $z_1$, whose integrand is
integrable by Lemma~\ref{lem:kg-linear}, it gives $G^+f(q)>0$. The function $G^+f$ is continuous,
nonnegative (its integrand is) and supported in the future of a compact set
(Proposition~\ref{prop:kg-ext}), so its integrand is integrable (Lemma~\ref{lem:kg-linear}). For
$h=G^+f$ at $p$ with $y_1$ the lemma gives $(G^+)^2f(p)>0$.
\end{proof}

\begin{lemma}[The tail]\label{lem:kg-tail}
Let $m^2,t_0,C\in\R$, and let $f$ be smooth with compact support, $f(q)=0$ for $q^0<t_0$, and
$|f|\le C$. For $p^0\ge t_0$ and $s=p^0-t_0$, the series of Definition~\ref{def:kg-series} satisfies
\[
\bigl|G^+_{m^2}f(p)-G^+f(p)+m^2(G^+)^2f(p)\bigr|\le C\,(m^2)^2\,s^6\cosh\bigl(\sqrt{|m^2|}\,s\bigr).
\]
\end{lemma}
\begin{proof}
The left side is $|\sum_{j\ge2}(-m^2)^j(G^+)^{j+1}f(p)|$. With $\xi=\sqrt{|m^2|}\,s$,
Proposition~\ref{prop:kg-conv} with $T=p^0$ bounds it by $Cs^2\sum_{j\ge2}\xi^{2j}/(2j)!$; here $C\ge0$,
since $|f|\le C$. And
$\sum_{j\ge2}\xi^{2j}/(2j)!=\xi^4\sum_{k\ge0}\xi^{2k}/(2k+4)!\le\xi^4\cosh\xi$, since $(2k+4)!\ge(2k)!$,
and $Cs^2\xi^4=C(m^2)^2s^6$.
\end{proof}

\begin{lemma}[Dilation]\label{lem:kg-dil}
Let $\lambda>0$, $f:M\to\R$ and $f_\lambda(q)=f(q/\lambda)$, with the extended Green operators, iterates and series of
Definitions~\ref{def:kg-supported} and~\ref{def:kg-series}.
\begin{enumerate}[label=(\alph*),leftmargin=2.5em]
\item For continuous $f$ and every $p$, $G^+f_\lambda(\lambda p)=\lambda^2\,G^+f(p)$.
\item For smooth $f$ with compact support and every $m^2\in\R$,
$(G^+)^jf_\lambda(q)=\lambda^{2j}(G^+)^jf(q/\lambda)$ for all $j$ and $q$, and
$G^+_{m^2}f_\lambda(\lambda p)=\lambda^2\,G^+_{\lambda^2m^2}f(p)$.
\end{enumerate}
\end{lemma}
\begin{proof}
\begin{enumerate}[leftmargin=2.5em]
\item[(a)] Substitute $y=\lambda z$: $dy=\lambda^3dz$ and $|y|=\lambda|z|$, which leaves $\lambda^2$. One
integral converges absolutely exactly when the other does.
\item[(b)] Induction on $j$, using (a) for the continuous function $(G^+)^jf$
(Corollary~\ref{cor:kg-iterate}) and homogeneity. Then
$\sum_j(-m^2)^j\lambda^{2j+2}(G^+)^{j+1}f(p)=\lambda^2\sum_j(-\lambda^2m^2)^j(G^+)^{j+1}f(p)$.
\end{enumerate}
\end{proof}

\begin{theorem}[Failure of Huygens' principle for $m^2\neq0$]\label{thm:kg-fail}
For every $m^2\neq0$, with the series of Definition~\ref{def:kg-series} and $\ll$ as in
Definition~\ref{def:kg-cone}, there are a test function $f$ and a point $p$ with $q\ll p$ for all
$q\in\supp f$ and $G^+_{m^2}f(p)\ne0$. Likewise there are $f\in\cD$ and $p$ with $p\ll q$ for all
$q\in\supp f$ and $G^-_{m^2}f(p)\neq0$.
\end{theorem}
\begin{proof}
\emph{A fixed configuration.} Let $f_0$ be a smooth bump function with $0\le f_0\le1$, $f_0(0)=1$ and
support in $\{\max(|q^0|,|\vec q|)\le\tfrac14\}$, and let $p_0=(1,\vec0)$. Every $q\in\supp f_0$ has
$|\vec p_0-\vec q|\le\tfrac14<\tfrac34\le p_0^0-q^0$, so $q\ll p_0$. Hence $G^+f_0(p_0)=0$
(Proposition~\ref{prop:kg-huygens}) and $c=(G^+)^2f_0(p_0)>0$ (Proposition~\ref{prop:kg-pos2} with
$q_0=0$).

\emph{Small parameters.} Lemma~\ref{lem:kg-tail} with $C=1$, $t_0=-\tfrac14$ and $s=\tfrac54$ gives
$|G^+_{m_0^2}f_0(p_0)+m_0^2c|\le(m_0^2)^2s^6\cosh s$ for $|m_0^2|\le1$, since $\cosh$ is increasing
on $[0,\infty)$. If moreover $m_0^2\ne0$ and $|m_0^2|\,s^6\cosh s<c$, the right side is less than
$|m_0^2|\,c$, so $G^+_{m_0^2}f_0(p_0)\ne0$.

\emph{Every $m^2$.} Choose $\lambda>0$ so small that $m_0^2=\lambda^2m^2$ satisfies both conditions,
and put $f=(f_0)_\lambda$, $p=\lambda p_0$. Then $\supp f=\lambda\supp f_0$, and $q\ll p$ if and only if
$q/\lambda\ll p_0$. By Lemma~\ref{lem:kg-dil}, $G^+_{m^2}f(p)=\lambda^2G^+_{m_0^2}f_0(p_0)\ne0$.

\emph{The advanced case.} Take $g=f\circ r$ and the point $rp$. Then
$G^-_{m^2}g(rp)=G^+_{m^2}f(p)\ne0$ by Proposition~\ref{prop:kg-adv}, and $q\ll p$ for all
$q\in\supp f$ becomes $rp\ll q'$ for all $q'\in\supp g$.
\end{proof}

\begin{theorem}[Huygens' principle holds exactly at $m^2=0$]\label{thm:kg-iff}
For $m^2\in\R$, the series $G^+_{m^2}$ of Definition~\ref{def:kg-series} and $\ll$ as in
Definition~\ref{def:kg-cone}, the following are equivalent:
\begin{enumerate}[label=(\roman*),leftmargin=2.5em]
\item $G^+_{m^2}f(p)=0$ for every $f\in\cD$ and every $p\in M$ with $q\ll p$ for all $q\in\supp f$;
\item $m^2=0$.
\end{enumerate}
The same holds for $G^-_{m^2}$, with $p\ll q$ in (i).
\end{theorem}
\begin{proof}
(ii)$\Rightarrow$(i) is Corollary~\ref{cor:kg-zero} with Proposition~\ref{prop:kg-huygens};
(i)$\Rightarrow$(ii) is Theorem~\ref{thm:kg-fail}.
\end{proof}
\noindent\motiv{Condition (i) is not vacuous: the configuration of
Theorem~\ref{thm:kg-fail} has $\supp f\neq\emptyset$ and $p$ in the open future cone of each of its
points.}

\medskip\noindent\motiv{For $m^2\neq0$ the Green data are not trivial: \revo{for the test function $f$ and the point $p$ of the second sentence of Theorem~\ref{thm:kg-fail} (the advanced case)}, $G^+_{m^2}f(p)=0$ because $p\notin J^+(\supp f)$, so
$Ef(p)=G^-_{m^2}f(p)\neq0$, and $\sigma\neq0$ by Proposition~\ref{prop:green-nondeg}.}

\medskip\noindent\motiv{By Theorem~\ref{thm:kg-field}(b), Huygens' principle is a property of $P_{m^2}$, not of the series.
For operators with the supports of Definition~\ref{def:kg-greenops}(b), it is equivalent to the
support form of Proposition~\ref{prop:kg-huygens}, on paper: if $p\notin\supp f+\partial C^+$, then
$\supp f$ lies in the union of the disjoint open sets $p-I^+$ and $M\setminus(p-C^+)$, and splitting
$f$ accordingly gives one part covered by (i) and one part whose support misses $J^-(\{p\})$.}

\begin{remark}[Kernels]\label{rem:kg-kernels}
Condition (i) of Theorem~\ref{thm:kg-iff} says that the convolution kernel of $G^+_{m^2}$, a
distribution in $(t,x)\in M$ (its Schwartz kernel is $k(p-q)$), vanishes on the open future cone
$I^+$ (on paper; kernels are not formalized). For $m^2=0$ the kernel is $\delta(t-|x|)/(4\pi|x|)$ (Remark~\ref{rem:spherical-means}),
and the kernel of $(G^+)^2$ is $1/(8\pi)$ times the indicator function of $C^+$, so that
$c=\frac1{8\pi}\int_{p_0-C^+}f_0$ in Theorem~\ref{thm:kg-fail}.

For $m^2>0$, with $m=\sqrt{m^2}$, the kernel has, on $I^+$, the additional term
\[
-\frac m{4\pi}\,\frac{J_1\bigl(m\sqrt{t^2-|x|^2}\bigr)}{\sqrt{t^2-|x|^2}},
\]
with $J_1$ the Bessel function~\vetsource{\cite{gunther} (pinpoint to be checked)}. For $m^2<0$ the additional term is
\revm{\[
\frac{\sqrt{-m^2}}{4\pi}\,\frac{I_1\bigl(\sqrt{-m^2}\sqrt{t^2-|x|^2}\bigr)}{\sqrt{t^2-|x|^2}},
\]}
with the modified Bessel function $I_1$. Summing the kernels of the powers $(G^+)^{j+1}$, which are locally integrable
functions for $j\ge1$, gives these terms.

The proof above uses neither formula, only the first two terms of the series and a bound on the
rest. For $m^2<0$ and $f\ge0$ every term $(-m^2)^j(G^+)^{j+1}f$ is nonnegative, so
Proposition~\ref{prop:kg-pos2} alone gives $G^+_{m^2}f(p)>0$ at every $p$ in the open future cone of
a point where $f>0$, without the tail bound or dilation. For $m^2>0$, $f\ge0$ and $p$ in the open
future cone of every point of $\supp f$, Lemma~\ref{lem:kg-tail} gives $G^+_{m^2}f(p)<0$ for the
dilated configurations of Theorem~\ref{thm:kg-fail} with small $\lambda$. The additional kernel term
changes sign at the zeros of $J_1(m\sqrt{t^2-|x|^2})$, and so does $G^+_{m^2}f(p)$ for $f$ concentrated
near a point (on paper). Bessel functions are not formalized.
\end{remark}
\end{extra}
\endgroup
